\section{Negative heat remains positive on every fixed unitary group}
\label{sec:analytic-finite-rank-characters}

Keep the normalization
\[
 p_1(t)=F_t'(0)>0,\qquad
 E_t(w)=F_t(w/p_1(t)),\qquad
 \Psi_t(z)=E_t(z-1)=\sum_{k\ge0}a_k(t)z^k.
\]
Thus $E_t(0)=E_t'(0)=1$ and $\Psi_t(1)=1$ for every real $t$.
For $U\in U(d)$ with eigenvalues $z_1,\ldots,z_d$, define
\[
             \chi_{t,d}(U)=\prod_{i=1}^d\Psi_t(z_i).
\]
Here a normalized character means a continuous central function of
positive type with value one at the identity. It need not be the
character of a finite-dimensional representation.

\begin{theorem}[NN-A021: positivity at fixed unitary rank]
Assume A. For every fixed $d\ge1$ there is $\eta_d>0$ such that
$\chi_{t,d}$ is a normalized character of $U(d)$ for
$-\eta_d\le t\le0$. Its expansion is
\begin{equation}\label{eq:analytic-fixed-unitary-expansion}
 \chi_{t,d}(U)=
 \sum_{\ell(\lambda)\le d}c_{\lambda,d}(t)s_\lambda(U),
 \qquad
 c_{\lambda,d}(t)=
       \det\bigl[a_{\lambda_i-i+j}(t)\bigr]_{i,j=1}^d>0,
\end{equation}
where $a_k=0$ for $k<0$ and partitions are padded with zero parts.
The series converges absolutely and uniformly on $U(d)$, and
\[
       \sum_{\ell(\lambda)\le d}
          c_{\lambda,d}(t)s_\lambda(1^d)=1.
\]
In particular every polynomial irreducible representation has
strictly positive probability weight, with no bound on its degree.

Under A and Rodgers--Tao, for every $t<0$ the compatible functions
$\chi_{t,d}$ fail positive type at some finite rank. If
\[
 d_*(t)=\min\{d\ge1:\chi_{t,d}\text{ is not of positive type}\},
\]
then $d_*(t)<\infty$ for $t<0$ and
\[
                         d_*(t)\longrightarrow\infty
                                  \qquad(t\uparrow0).
\]
For $t\ge0$, all ranks have positive type by NN-A008. No
simplicity assumption is needed at $t=0$.
\end{theorem}

\begin{proof}
We first establish positivity of the Taylor coefficients in the
required Toeplitz sense. The order of $E_t$ is less than one, and its
genus-zero product is
\[
 E_t(w)=\prod_j(1+\beta_j(t)w),\qquad
 \sum_j|\beta_j(t)|<\infty,\qquad
 \sum_j\beta_j(t)=1.
\]
The factors are real positive or occur in conjugate pairs when
$t$ is sufficiently close to zero, with angles tending uniformly
to zero by NN-A020. More explicitly, a zero $x$ of $H_t$ in the
right half-plane gives $\beta=1/(p_1(t)x^2)$, so that
$|\arg\beta|=2|\arg x|$ near zero.

Under A, all $\beta_j(0)$ are positive. There are infinitely many,
and their sum is one, hence $\beta_{\max}(0)<1$. Choose
$b$ with $\beta_{\max}(0)<b<1$. The entire function $E_0$ is
zero-free on $|w|\le1/b$. Local uniform convergence
$E_t\to E_0$ makes this disc zero-free for all sufficiently small
$|t|$. Consequently
\begin{equation}\label{eq:analytic-unitary-beta-bound}
                |\beta_j(t)|<b\qquad\text{for every }j.
\end{equation}
This argument does not require individually labelling any multiple
zero. Define
\[
       \widetilde\beta_j(t)=\frac{\beta_j(t)}{1-\beta_j(t)}.
\]
For $\beta=\rho e^{i\theta}$, $\rho\le b<1$, with $|\theta|<\pi/2$,
\[
 |\arg\widetilde\beta|
 \le |\theta|+
       \arctan\frac{\rho|\sin\theta|}{1-\rho\cos\theta}
 \le \frac{|\theta|}{1-b}.
\]
Thus the angles of all $\widetilde\beta_j$ tend uniformly to zero,
and their absolute values are summable.

The real factors $1-\beta_j$ are positive; each conjugate pair
has product $|1-\beta_j|^2>0$. Absolute convergence therefore gives
\begin{equation}\label{eq:analytic-shifted-PF-product}
 \Psi_t(z)=E_t(-1)\prod_j(1+\widetilde\beta_j(t)z),
                  \qquad E_t(-1)>0.
\end{equation}
For a fixed $d$, make $|t|$ small enough that every conjugate-pair
angle is less than $\pi/(d+1)$. The quadratic factor lemma in
NN-A020 proves $PF_d$ for the coefficients of the normalized
product in \eqref{eq:analytic-shifted-PF-product}. Positive
constant scaling preserves every minor sign. Moreover every
allowed minor of order at most $d$ is strictly positive: the
infinite product has infinitely many factors, all with positive
linear coefficient, and the one-step particle paths in the
Cauchy--Binet proof of NN-A020 apply unchanged. This establishes
strict positivity at arbitrary indices, not merely in a bounded
coefficient window. Shrinking the neighbourhood gives a closed
interval $[-\eta_d,0]$.

We give the character expansion explicitly. For any entire
function $\Psi(z)=\sum_{k\ge0}a_kz^k$, multiply
$\prod_i\Psi(x_i)$ by
$\Delta(x)=\det[x_i^{d-j}]_{i,j=1}^d$ and expand the rows:
\[
 \det\bigl[x_i^{d-j}\Psi(x_i)\bigr]_{i,j=1}^d.
\]
Cauchy--Binet, with the exponents chosen in strictly decreasing
order $m_i=\lambda_i+d-i$, gives the coefficient of the alternant
$\det[x_i^{\lambda_j+d-j}]$ as
$\det[a_{\lambda_i-i+j}]$. Dividing by $\Delta$ gives
\begin{equation}\label{eq:analytic-alternant-formal}
 \prod_{i=1}^d\Psi(x_i)
     =\sum_{\ell(\lambda)\le d}
        \det[a_{\lambda_i-i+j}]s_\lambda(x_1,\ldots,x_d)
\end{equation}
as a degree-wise formal identity. Each degree involves only
finitely many terms. The Schur polynomials are the irreducible
polynomial characters by the Weyl character formula.

For $\Psi=\Psi_t$, the coefficient determinant is the Toeplitz
minor with row indices $i-\lambda_i$ and column indices $j$,
after adding the same sufficiently large integer to both lists.
The row indices increase strictly, and
$i-\lambda_i\le i$, so this minor is allowed. It is therefore
strictly positive. Set every $x_i=u$, $0\le u<1$, in
\eqref{eq:analytic-alternant-formal}. Nonnegative coefficients
and degree-wise convergence give
\[
 \Psi_t(u)^d
      =\sum_{\ell(\lambda)\le d}
        c_{\lambda,d}(t)u^{|\lambda|}s_\lambda(1^d).
\]
Monotone convergence as $u\uparrow1$ proves that the positive
weights $c_{\lambda,d}(t)s_\lambda(1^d)$ sum to one.
For $U\in U(d)$,
$|s_\lambda(U)|\le s_\lambda(1^d)$, since this is the trace of
a unitary representation. The series thus converges absolutely
and uniformly. The formal identity inside the unit polydisc,
followed by radial convergence, identifies its value with
$\prod_i\Psi_t(z_i)$. A positive mixture of normalized unitary
characters is of positive type, proving the first assertion.

For completeness we prove the required all-ranks converse,
including a possible zero constant coefficient, rather than
postulating an infinite-group classification. Suppose an entire
function $\Psi$ with real Taylor coefficients and $\Psi(1)=1$
has $\prod_i\Psi(z_i)$ of positive type on every $U(d)$.
At rank one its Fourier coefficients are nonnegative; these
are exactly its Taylor coefficients $a_k$, with the negative
coefficients zero. Let $m$ be the first index with $a_m>0$ and
write $\Psi(z)=z^m B(z)$. Then $B(0)>0$, $B(1)=1$, and
every product $\prod_i B(z_i)$ is again of positive type:
it is the original function multiplied by the one-dimensional
unitary character $\det(U)^{-m}$.

The coefficient of the irreducible polynomial character
$s_\lambda$ in this central function is
\begin{equation}\label{eq:analytic-necessary-Schur-coefficient}
                      \det[b_{\lambda_i-i+j}]_{i,j=1}^d,
                       \qquad B(z)=\sum_{k\ge0}b_kz^k.
\end{equation}
One may obtain this identity directly from the Weyl integration
formula: multiply the two alternants, expand their determinants,
and use
$\int_{S^1}z^n\,dz/(2\pi iz)=\mathbf1_{n=0}$.
Since a continuous central positive-type function on a compact
group has nonnegative irreducible-character coefficients,
\eqref{eq:analytic-necessary-Schur-coefficient} is nonnegative.
This last fact also follows directly from a positive-type
quadratic form integrated against matrix coefficients, followed
by Schur orthogonality. Thus the specialization
\[
             \rho(h_k)=b_k/b_0,\qquad h_0=1,
\]
is nonnegative on every ordinary Schur function, because every
partition occurs at some finite rank.

Here ordinary Schur positivity implies skew Schur positivity.
We recall a finite-dimensional proof of the needed implication.
In two disjoint alphabets $X,Y$, splitting a semistandard
tableau according to its alphabet gives
\[
        s_\lambda(X,Y)=
             \sum_{\mu\subseteq\lambda}s_\mu(X)s_{\lambda/\mu}(Y).
\]
On the other hand, restrict the polynomial irreducible
representation with character $s_\lambda$ of $U(p+q)$ to
$U(p)\times U(q)$, where $p,q$ exceed the diagram size.
Complete reducibility gives a nonnegative integer sum
\[
        s_\lambda(X,Y)=
             \sum_{\mu,\nu}N_{\mu,\nu}^\lambda
                           s_\mu(X)s_\nu(Y),
                         \qquad N_{\mu,\nu}^\lambda\ge0.
\]
The summands are polynomial representations, so only partitions
occur. Linear independence of the $s_\mu(X)$ identifies
$s_{\lambda/\mu}(Y)=\sum_\nu N_{\mu,\nu}^\lambda s_\nu(Y)$.
This is an identity of symmetric functions by the choice of
large alphabets. Applying $\rho$ proves nonnegativity of every
skew Schur function. The ordinary Jacobi--Trudi formula
identifies every allowed Toeplitz minor of $(b_k/b_0)$ with
one of these values; forbidden minors are zero. Thus $(b_k)$
is $PF_\infty$.

Edrei's theorem now says that the entire generating function
$B$ has only nonpositive real zeros: its possible denominator
factors would have positive real poles and cannot cancel its
nonpositive numerator zeros. The factor $z^m$ adds only the
zero at zero. We have proved that an entire $\Psi$ of the stated
all-ranks positive type cannot have a nonreal zero.

Apply this conclusion to $\Psi_t$. For every $t<0$, A and
Rodgers--Tao give a nonreal zero $x$ of $H_t$. There is no
zero on the imaginary axis since $H_t(iy)>0$. Hence
$x^2$ is nonreal, and
\[
                        z=1-p_1(t)x^2
\]
is a nonreal zero of $\Psi_t$. All-ranks positive type is
therefore impossible, so $d_*(t)$ is finite. Compatibility
\[
        \chi_{t,d+1}(U\oplus1)=\chi_{t,d}(U)
\]
shows that positivity at rank $d+1$ implies positivity at
rank $d$. The first part of the theorem then proves
$d_*(t)\to\infty$ as $t\uparrow0$.
\end{proof}

\begin{corollary}[The exact finite-rank dictionary]
Let $\Psi(z)=\sum_{k\ge0}a_kz^k$ be entire with real coefficients,
$a_0>0$ and $\Psi(1)=1$. For a fixed $d\ge1$,
\[
 \bigl[\,U\longmapsto\det\Psi(U)
                 \text{ is of positive type on }U(d)\,\bigr]
 \quad\Longleftrightarrow\quad
                  (a_k)_{k\ge0}\in PF_d.
\]
Consequently, for $t<0$ sufficiently close to zero, $d_*(t)$ is
exactly the minimum order of a negative Toeplitz minor of the
coefficients of $\Psi_t$. These are the coefficients of the
translated normalized function $\Psi_t(z)=E_t(z-1)$, not the
untranslated coefficients of $F_t$.
\end{corollary}

\begin{proof}
The forward construction from $PF_d$ to a positive normalized
Schur expansion was proved above and did not require strict
inequalities. Conversely positive type gives the nonnegative
coefficient determinants
\eqref{eq:analytic-necessary-Schur-coefficient} for every
partition of length at most $d$. Normalize $a_0$ to one in the
symmetric-function specialization. In the preceding skew Schur
expansion, the coefficient $N_{\mu,\nu}^\lambda$ vanishes unless
$\mu\subseteq\lambda$, by the two-alphabet tableau identity.
Taking equally large alphabets and interchanging their two
unitary blocks proves
$N_{\mu,\nu}^\lambda=N_{\nu,\mu}^\lambda$; hence it also
vanishes unless $\nu\subseteq\lambda$.

A Toeplitz minor of order $k\le d$, with increasing indices
$i_1,\ldots,i_k$ and $j_1,\ldots,j_k$, is zero when some
$i_\ell>j_\ell$. Otherwise transpose and reverse the index
lists, and put
\[
 \lambda_p=j_{k+1-p}+p-i_1-k,\qquad
 \mu_p=i_{k+1-p}+p-i_1-k,\qquad 1\le p\le k.
\]
These are partitions with $\mu\subseteq\lambda$ and
$\ell(\lambda)\le k$. The normalized determinant is
$\rho(s_{\lambda/\mu})$ by the ordinary Jacobi--Trudi identity.
Its expansion uses only shapes $\nu\subseteq\lambda$, so every
$\ell(\nu)\le k\le d$. All their specialized Schur values are
nonnegative by the character coefficients already obtained.
Thus every minor of order at most $d$ is nonnegative.
The minimum-order assertion follows from the equivalence and
$E_t(-1)>0$ near zero.
\end{proof}

\paragraph{Effectivity at fixed rank.}
The interval $\eta_d$ can be certified effectively under A.
First certify that $E_0$ is zero-free on a rational-radius disc
$|w|\le1/b$ for some rational $b<1$. Such a choice exists by
$\beta_{\max}(0)<1$, and rational grid evaluation with a
derivative bound certifies a positive compact minimum.
Uniform effective bounds for $\partial_t E_t$ on a slightly
larger compact set, obtained from the defining theta integrals
and the positive denominator $p_1(t)$, then give a rational
time interval on which \eqref{eq:analytic-unitary-beta-bound}
holds. In the effective corollary to NN-A020 choose its integer
parameter $r$ large enough that
\[
       \frac{2\arctan(1/[4(r+1)])}{1-b}
                                      <\frac{\pi}{d+1}.
\]
That corollary certifies exactly the angular confinement used
above. The smaller of the two time intervals is a certified
$\eta_d$. This asserts a terminating finite computation at
each fixed rank; it does not assert an asymptotic rate or
polynomial running time.

\paragraph{Primary character locators and geometric translation.}
The Weyl character formula used here is, for example, Etingof
et al., \emph{Introduction to Representation Theory},
Theorem~5.22.1, printed p.~132; the neighbouring
Theorem~5.23.2 gives complete reducibility of algebraic
representations. For the precise compact-group convention,
Faraut's \emph{Infinite Dimensional Spherical Analysis},
notes written by Sho Matsumoto, Kyushu (September 2007), section~4.1,
printed p.~56, states the positive normalized-character
expansion, and section~4.2, printed p.~57, gives the Weyl
integration and character formulas. The new conclusion is
the application to bounded negative heat time at every fixed
rank, not those classical character formulas.

Exterior powers and their additive filtration formula have
the concrete Grothendieck locator SGA~6, Expos\'e~V,
Example~2.2(b) and Lemma~2.2.1, including the representation-ring
case. The coefficients in
\eqref{eq:analytic-fixed-unitary-expansion} give positive
probability measures on actual irreducible polynomial
$U(d)$-representations. They are real weights, not asserted
integer multiplicities in an algebraic vector bundle.
The analytic heat path lies in each fixed-rank positive
character cone on some negative interval, while it lies
outside the compatible positive cone for all ranks at every
negative time. This is an unbounded-degree statement at
fixed rank, stronger than finitely many character tests;
it neither supplies an integral positive basis nor a GCT
obstruction or a lower bound against all algorithms.
